\chapter{Asepsie et antisepsie}

\noindent\textit{Après avoir vu comment les microorganismes pénètrent dans
l'organisme, ce chapitre présente les moyens dont nous disposons pour les empêcher
d'y entrer, ou pour les éliminer lorsqu'ils sont déjà présents : l'asepsie,
l'antisepsie, et l'usage du préservatif.}
\vspace{8pt}

\connecteBanner

\section{L'asepsie : empêcher toute contamination}

\begin{definition}[Asepsie]
L'\textbf{asepsie} est l'ensemble des méthodes qui \textbf{empêchent} l'introduction
de microorganismes dans un lieu, un objet ou un organisme initialement exempt de
microorganismes pathogènes. C'est une démarche \textbf{préventive}.
\end{definition}

\begin{propriete}[Exemples de mesures d'asepsie]
\begin{itemize}
\item la \textbf{stérilisation} du matériel médical (par la chaleur, à l'autoclave) ;
\item l'utilisation de matériel \textbf{à usage unique} (aiguilles, seringues,
gants) ;
\item le \textbf{lavage des mains} avant un soin ;
\item le port d'un \textbf{masque} et de \textbf{gants} par le personnel soignant.
\end{itemize}
\end{propriete}

\begin{exemple}
Avant une opération chirurgicale, tous les instruments sont stérilisés à l'autoclave :
c'est une mesure d'asepsie, qui garantit qu'aucun microorganisme n'est introduit dans
le corps du patient pendant l'intervention.
\end{exemple}

\section{L'antisepsie : éliminer les microorganismes déjà présents}

\begin{definition}[Antisepsie]
L'\textbf{antisepsie} est l'application d'un produit, l'\textbf{antiseptique}, qui
détruit ou inhibe les microorganismes déjà présents sur la peau ou une plaie. C'est
une démarche \textbf{curative}, appliquée après un contact déjà survenu.
\end{definition}

\begin{propriete}[Exemples d'antiseptiques]
Alcool, eau de Javel diluée, solution iodée (type Bétadine), eau oxygénée : ces
produits sont appliqués directement sur la peau, une plaie, ou une surface, pour
réduire le nombre de microorganismes présents.
\end{propriete}

\begin{illus}
\begin{tikzpicture}[scale=0.85]
\node[font=\bfseries\small] at (0,1.5) {Asepsie};
\node[font=\scriptsize,align=center] at (0,0.6) {\textit{avant} le contact};
\draw[onbo,fill=onbo,fill opacity=0.12,line width=1.2pt] (0,-0.5) circle (0.75);
\node[font=\scriptsize,align=center] at (0,-0.5) {aucun\\microbe};
\node[font=\bfseries\small] at (4.5,1.5) {Antisepsie};
\node[font=\scriptsize,align=center] at (4.5,0.6) {\textit{après} le contact};
\draw[onbink,fill=onbink,fill opacity=0.12,line width=1.2pt] (4.5,-0.5) circle (0.75);
\node[font=\scriptsize,align=center] at (4.5,-0.5) {microbes\\détruits};
\end{tikzpicture}
\fig{Figure — Asepsie (prévention, avant contact) et antisepsie (élimination, après
contact).}
\end{illus}

\begin{exemple}
Désinfecter une plaie avec de l'alcool après une chute est une mesure d'antisepsie :
le microorganisme a déjà pu entrer en contact avec la plaie, et le produit vise à
l'éliminer.
\end{exemple}

\begin{remarque}
Piège fréquent : asepsie et antisepsie ne s'opposent pas, elles se complètent. Au bloc
opératoire, l'asepsie du matériel est associée à l'antisepsie de la peau du patient
(désinfection de la zone à opérer) avant l'intervention.
\end{remarque}

\section{Le préservatif, une barrière contre la contamination}

\begin{propriete}[Un moyen de prévention barrière]
Le \textbf{préservatif} constitue une \textbf{barrière physique} qui empêche le
contact direct entre les sécrétions sexuelles de deux partenaires. Il réduit ainsi le
risque de contamination par voie sexuelle (VIH, autres infections sexuellement
transmissibles) et joue également un rôle contraceptif.
\end{propriete}

\begin{remarque}
Le préservatif se rapproche de la logique de l'asepsie : il agit de façon
\textbf{préventive}, en empêchant le contact avec un microorganisme éventuellement
présent, plutôt qu'en le détruisant après coup.
\end{remarque}

% ═══════════════════════════════════════════════════════════════
\section{L'essentiel à retenir}

\begin{synthese}
\textbf{Asepsie.} Ensemble de méthodes préventives empêchant l'introduction de
microorganismes (stérilisation, matériel à usage unique, lavage des mains).\\[4pt]
\textbf{Antisepsie.} Application d'un antiseptique qui détruit les microorganismes
déjà présents sur la peau ou une plaie.\\[4pt]
\textbf{Antiseptiques courants.} Alcool, eau de Javel diluée, solution iodée, eau
oxygénée.\\[4pt]
\textbf{Préservatif.} Barrière physique préventive contre la contamination par voie
sexuelle.\\[4pt]
\textbf{Piège fréquent.} Asepsie (avant contact) et antisepsie (après contact) sont
complémentaires, non opposées.
\end{synthese}

% ═══════════════════════════════════════════════════════════════
\section{Méthodes \& savoir-faire}

\methode{Distinguer une mesure d'asepsie d'une mesure d'antisepsie}
Se demander si la mesure agit \textbf{avant} tout contact avec un microorganisme
(asepsie) ou \textbf{après} un contact déjà survenu, pour l'éliminer (antisepsie).
\begin{exemple}
Stériliser une seringue avant usage : asepsie. Désinfecter une plaie déjà présente :
antisepsie.
\end{exemple}

% ═══════════════════════════════════════════════════════════════
\section{Exercices d'application}
\begin{multicols}{2}

\rubrique{Asepsie}
\exo{}\diff{1} Définir l'asepsie.

\exo{}\diff{2} Citer deux mesures d'asepsie utilisées en milieu médical.

\rubrique{Antisepsie}
\exo{}\diff{1} Définir l'antisepsie.

\exo{}\diff{2} Citer deux exemples d'antiseptiques.

\rubrique{Asepsie et antisepsie}
\exo{}\diff{2} Quelle est la principale différence entre asepsie et antisepsie ?

\rubrique{Le préservatif}
\exo{}\diff{1} Contre quel type de contamination le préservatif protège-t-il ?

% ═══════════════════════════════════════════════════════════════
\end{multicols}

\section{Exercices d'entraînement}
\begin{multicols}{2}

\rubrique{Asepsie et antisepsie}
\exo{}\diff{2} Classer chacune des mesures suivantes en asepsie ou antisepsie :
stérilisation d'un scalpel, application d'alcool sur une plaie, port de gants,
lavage d'une coupure à l'eau de Javel diluée.

\exo{}\diff{3} Expliquer pourquoi, avant une opération chirurgicale, on utilise à la
fois des mesures d'asepsie et des mesures d'antisepsie.

\exo{}\diff{2} Un élève affirme que l'asepsie et l'antisepsie sont deux mots pour la
même chose. Corriger cette affirmation.

\rubrique{Le préservatif}
\exo{}\diff{3} Expliquer pourquoi le préservatif est comparé à une mesure d'asepsie
plutôt qu'à une mesure d'antisepsie.

% ═══════════════════════════════════════════════════════════════
\end{multicols}

\section{Sujets type BEPC}

\sujetbac[Asepsie et antisepsie]
\begin{enumerate}[label=\arabic*)]
\item Définis l'asepsie et l'antisepsie.
\item Cite deux mesures d'asepsie et deux exemples d'antiseptiques.
\item Un infirmier stérilise ses instruments puis désinfecte la peau du patient avec
de l'alcool avant une injection. Identifie, pour chaque étape, s'il s'agit d'asepsie
ou d'antisepsie.
\item Pourquoi ces deux démarches sont-elles complémentaires ?
\end{enumerate}

% ═══════════════════════════════════════════════════════════════
\section{Corrigés}
\begin{multicols}{2}

\corrige{de l'exercice 1}
L'ensemble des méthodes préventives qui empêchent l'introduction de microorganismes.

\corrige{de l'exercice 2}
Par exemple : la stérilisation du matériel et le lavage des mains.

\corrige{de l'exercice 3}
L'application d'un produit qui détruit ou inhibe les microorganismes déjà présents.

\corrige{de l'exercice 4}
Par exemple : l'alcool et l'eau de Javel diluée.

\corrige{de l'exercice 5}
L'asepsie agit avant tout contact avec un microorganisme (prévention) ; l'antisepsie
agit après un contact déjà survenu (élimination).

\corrige{de l'exercice 6}
Contre la contamination par voie sexuelle.

\corrige{de l'exercice 7}
Asepsie : stérilisation d'un scalpel, port de gants. Antisepsie : application d'alcool
sur une plaie, lavage d'une coupure à l'eau de Javel diluée.

\corrige{de l'exercice 8}
Parce que l'asepsie du matériel évite d'introduire de nouveaux microorganismes,
tandis que l'antisepsie de la peau élimine ceux déjà présents sur le patient : les
deux mesures réduisent ensemble le risque d'infection.

\corrige{de l'exercice 9}
C'est faux : l'asepsie prévient l'introduction de microorganismes avant tout contact,
alors que l'antisepsie élimine des microorganismes déjà présents ; ce sont deux
démarches complémentaires mais différentes.

\corrige{de l'exercice 10}
Parce qu'il agit en empêchant le contact avec un microorganisme éventuellement
présent, de façon préventive, plutôt qu'en détruisant un microorganisme déjà en
contact avec l'organisme.

\corrige{du sujet 1}
1) L'asepsie empêche préventivement l'introduction de microorganismes ; l'antisepsie
détruit les microorganismes déjà présents.\\
2) Asepsie : stérilisation, lavage des mains. Antiseptiques : alcool, eau de Javel
diluée.\\
3) Stériliser les instruments : asepsie. Désinfecter la peau à l'alcool : antisepsie.\\
4) Parce que l'une empêche l'introduction de nouveaux microorganismes et l'autre
élimine ceux déjà présents : ensemble, elles réduisent au maximum le risque
d'infection.

\end{multicols}
\exercicesBanner
